\section*{Chapter \ref{ch:b2:rigid-body-mechanics} --- Rigid-Body Mechanics}

\begin{solution}{exo:b2:rigid-body-mechanics:1}
(a) $\tfrac1{12}M\ell^2 = \qty{4.2e-2}{kg.m^2}$; through one end (Huygens,
$d = \ell/2$): $\tfrac13M\ell^2 = \qty{0.167}{kg.m^2}$. (b) $\tfrac12MR^2 =
\qty{4.0e-2}{kg.m^2}$; at the rim $\tfrac32MR^2 = \qty{0.12}{kg.m^2}$.
(c) $\tfrac25 \times 0.16 \times (0.0285)^2 = \qty{5.2e-5}{kg.m^2}$. (d) The
hoop, $MR^2$ against $\tfrac12MR^2$: twice.
\end{solution}

\begin{solution}{exo:b2:rigid-body-mechanics:2}
$v = \qty{5.0}{m/s}$, $\omega = v/R = 5/0.35 = \qty{14.3}{rad/s}$. Top:
$2v = \qty{10}{m/s}$ forward; bottom: $0$; front: $\vect v_G + \vect\Omega
\wedge\vect{GM} = (v, -v)$, i.e.\ $v\sqrt2 = \qty{7.1}{m/s}$ at
\qty{45}{\degree} downward. Relative to the bicycle the tread moves at
$R\omega = v = \qty{5}{m/s}$ everywhere. The stone leaves at \qty{10}{m/s},
horizontally forward.
\end{solution}

\begin{solution}{exo:b2:rigid-body-mechanics:3}
$J = \tfrac12 \times 40 \times 0.25^2 = \qty{1.25}{kg.m^2}$; $\omega = \qty{628}{rad/s}$.
(a) $E = \tfrac12J\omega^2 = \qty{246}{kJ}$. (b) $\Gamma = J\omega/\Delta t =
1.25 \times 628/120 = \qty{6.5}{N.m}$; $P = \Gamma\omega = \qty{4.1}{kW}$ at the end.
(c) \qty{200}{kJ} taken out leave \qty{46}{kJ}: $\omega = \sqrt{2 \times 46\,000/1.25}
= \qty{271}{rad/s} = \qty{2600}{rpm}$. (d) Air drag on the rim grows as
$\omega^2$ (power as $\omega^3$) and would heat and slow the rotor; in
vacuum on magnetic \omterm{def:b2:rigid-body-mechanics:pivot}{bearings} the losses fall to a few watts.
\end{solution}

\begin{solution}{exo:b2:rigid-body-mechanics:4}
$J = \tfrac12MR^2 + MR^2 = \tfrac32MR^2$, $d = R$: $\ell_{\text{eq}} = \tfrac32R
= \qty{0.30}{m}$, $T = 2\pi\sqrt{0.30/9.81} = \qty{1.10}{s}$. At $R/2$: $J =
\tfrac12MR^2 + \tfrac14MR^2 = \tfrac34MR^2$, $d = R/2$, $\ell_{\text{eq}} =
\tfrac32R$ again, same period. In general $\ell_{\text{eq}}(d) = (R^2/2 +
d^2)/d$, minimal at $d = R/\sqrt2$: $\ell_{\text{eq}} = R\sqrt2 = \qty{0.283}{m}$,
$T = \qty{1.07}{s}$.
\end{solution}

\begin{solution}{exo:b2:rigid-body-mechanics:5}
(a) $\tan20^\circ = 0.36 \le f(1 + k)/k$, i.e.\ $\le 0.6$ (hoop), $0.9$
(disk), $1.05$ (sphere): all roll. (b) $v = \sqrt{2gh/(1 + k)}$: hoop
\qty{4.4}{m/s}, disk \qty{5.1}{m/s}, sphere \qty{5.3}{m/s}; sphere first, hoop
last. (c) Length $h/\sin\alpha = \qty{5.85}{m}$, $a = g\sin\alpha/(1 + k) =
1.68$, $2.24$, $\qty{2.40}{m/s^2}$, $t = \sqrt{2d/a} = 2.6$, $2.3$, \qty{2.2}{s}.
(d) $T = \dfrac{k}{1 + k}Mg\sin\alpha = \tfrac27Mg\sin\alpha$; the contact
point is at rest, so $\vect T\cdot\vect v_I = 0$. (e) Hoop: $\tan\alpha >
2f = 0.6$, $\alpha > \qty{31}{\degree}$; sphere: $\tan\alpha > 3.5f = 1.05$,
$\alpha > \qty{46}{\degree}$.
\end{solution}

\begin{solution}{exo:b2:rigid-body-mechanics:6}
$\Gamma = 2fNr = 2 \times 0.4 \times 800 \times 0.12 = \qty{77}{N.m}$; $\dot\omega =
\Gamma/J = \qty{64}{rad/s^2}$; $\omega_0 = \qty{94}{rad/s}$, $t = \qty{1.5}{s}$;
$\theta = \omega_0^2/2\dot\omega = \qty{69}{rad} = 11$ turns; heat $= \tfrac12J
\omega_0^2 = \qty{5.3}{kJ}$, $\Delta T = 5300/(1.5 \times 470) = \qty{7.6}{K}$.
\end{solution}

\begin{solution}{exo:b2:rigid-body-mechanics:7}
(a) The rope must exert a net torque to spin the pulley up: $T_1 \ne T_2$.
(b) $m_1a = m_1g - T_1$, $m_2a = T_2 - m_2g$, $J\dot\omega = (T_1 - T_2)R$ with
$a = R\dot\omega$: $a = (m_1 - m_2)g/(m_1 + m_2 + J/R^2) = 0.5 \times 9.81/4.0 =
\qty{1.23}{m/s^2}$. (c) $T_1 = m_1(g - a) = \qty{17.2}{N}$, $T_2 = m_2(g + a) =
\qty{16.6}{N}$. (d) Massless: \qty{1.40}{m/s^2}. Energy: $(m_1 - m_2)gx =
\tfrac12(m_1 + m_2 + J/R^2)v^2$; differentiate.
\end{solution}

\begin{solution}{exo:b2:rigid-body-mechanics:8}
(a) Weight $mg$ at the middle; wall: horizontal $N_w$ at the top; ground:
vertical $N = mg$ and horizontal friction $T = N_w$. Moments about the
foot: $mg\,\tfrac{\ell}2\cos\theta = N_w\ell\sin\theta$. (b) $T = N_w =
mg/(2\tan\theta) \le fmg \iff \tan\theta \ge 1/2f$. (c) $\tan\theta \ge
1.25$, $\theta \ge \qty{51}{\degree}$. (d) Person at $x$: $N = 4mg$, $N_w =
mg(\tfrac12 + 3x/\ell)/\tan60^\circ \le 4fmg = 1.6mg$: $x/\ell \le (2.77 -
0.5)/3 = 0.76$.
\end{solution}

\begin{solution}{exo:b2:rigid-body-mechanics:9}
(a) The contact point moves forward at $v + R\omega$ ($\omega > 0$ for
back-spin): friction $-fmg$ on the centre, $\dot v = -fg$; torque
$-fmgR$ on the spin: $\dot\omega = -5fg/2R$. (b) Slip $v + R\omega = v_0 +
R\omega_0 - \tfrac72fgt$ vanishes at $t_1 = 2(v_0 + R\omega_0)/7fg$, then $v_1
= v_0 - fgt_1 = (5v_0 - 2R\omega_0)/7$. (c) $v_1 < 0 \iff R\omega_0 > \tfrac52v_0$.
(d) $t_1 = 16/(7 \times 1.96) = \qty{1.2}{s}$; $v_1 = (10 - 12)/7 = \qty{-0.29}{m/s}$:
it comes back.
\end{solution}

\begin{solution}{exo:b2:rigid-body-mechanics:10}
(a) With $x$ toward the pull, friction $T$ along $x$: $Ma = F + T$;
moments about $G$ (counterclockwise): $rF + RT = -kMRa$ (rolling, $\omega_z
= -a/R$). Hence $a = F(R - r)/(1 + k)MR > 0$: the spool rolls
\emph{toward} the pull, winding the thread up. (b) $T = -F(kR + r)/(1 +
k)R$ (backward); no slip if $F \le fMg(1 + k)R/(kR + r)$. (c) The thread's
line of action is tangent to the hub; its moment about the contact
point $I$ is $F(r - R\cos\beta)$, clockwise (forward roll) when $\cos\beta >
r/R$, zero when the line passes through $I$ --- the instantaneous axis
--- and then $a = F(R\cos\beta - r)/(1 + k)MR$, $T = -F(kR\cos\beta + r)/(1 + k)R$,
$N = Mg - F\sin\beta$. (d) For $\cos\beta < r/R$ the moment about $I$ is
counterclockwise: the spool rolls \emph{away} from the puller.
\end{solution}

\begin{solution}{exo:b2:rigid-body-mechanics:11}
(a) $Mv_0 = P$; about $G$: $\tfrac25MR^2\omega_0 = P(h - R)$ (top-spin sense
for $h > R$). (b) $v_0 = R\omega_0 \iff 2R = 5(h - R) \iff h = \tfrac75R$; above
it $R\omega_0 > v_0$ (top-spin), below it back-spin or under-spin. (c)
$\omega_0 = 0$: slides for $t_1 = 2v_0/7fg = \qty{0.44}{s}$, distance $v_0t_1 -
\tfrac12fgt_1^2 = \qty{1.1}{m}$. (d) A cushion at $\tfrac75R$ returns the ball
\omterm{def:b2:rigid-body-mechanics:rolling}{rolling without slipping}: no sliding phase, no speed lost to the cloth,
a predictable rebound.
\end{solution}

\begin{solution}{exo:b2:rigid-body-mechanics:12}
(a) $v_G = (R - r)\dot\theta$; the contact point is at rest, so the
cylinder spins at $\omega = v_G/r = (R - r)\dot\theta/r$ about its axis.
(b) $E_k = \tfrac12Mv_G^2 + \tfrac12\cdot\tfrac12Mr^2\omega^2 = \tfrac34M(R - r)^2
\dot\theta^2$, $E_p = -Mg(R - r)\cos\theta$; $\dd E/\dd t = 0$ gives $\ddot\theta
+ \dfrac{2g}{3(R - r)}\sin\theta = 0$, $T = 2\pi\sqrt{3(R - r)/2g}$. (c) A
sliding point: $2\pi\sqrt{(R - r)/g}$ --- rolling is slower by $\sqrt{3/2}$,
a third of the energy being rotation. (d) Tangential: $M(R - r)\ddot\theta =
-Mg\sin\theta + T$, so $T = \tfrac13Mg\sin\theta$; at the turning point
$N = Mg\cos\theta_0$: $f \ge \tfrac13\tan\theta_0 \approx \theta_0/3$.
\end{solution}

\begin{solution}{pb:b2:rigid-body-mechanics:1}
\textbf{1.} $v = \qty{25}{m/s}$, $\omega = \qty{80.6}{rad/s}$; top \qty{50}{m/s},
bottom $0$, front $(25, -25)$: \qty{35}{m/s} at \qty{45}{\degree} downward.

\textbf{2.} $\tfrac12Mv^2 = \qty{375}{kJ}$; $4 \times \tfrac12J\omega^2 = \qty{13}{kJ}$;
total \qty{388}{kJ}, wheels $3.3\%$.

\textbf{3.} $N_f + N_r = Mg = \qty{11.8}{kN}$, $N_fa = N_rb$: $N_f = Mgb/L =
\qty{6.8}{kN}$, $N_r = Mga/L = \qty{5.0}{kN}$ --- the front (engine) axle.

\textbf{4.} $P_r = \sum\mu_rN_iR\omega = \mu_rMgv = 0.012 \times 11772 \times 25 =
\qty{3.5}{kW}$; drag $\tfrac12 \times 1.2 \times 0.7 \times 625 = \qty{262}{N}$, \qty{6.6}{kW}.

\textbf{5.} $(M + 4J/R^2)\dot v = -\mu_rMg - \tfrac12\rho C_xSv^2$, $M_{\text{eff}} =
1200 + 41.6 = \qty{1242}{kg}$. Rolling only: $\dot v = \qty{-0.114}{m/s^2}$,
$d = (25^2 - 12.5^2)/0.228 = \qty{2.1}{km}$. Drag only: $v = v_0\eu^{-x/\lambda}$,
$\lambda = M_{\text{eff}}/\tfrac12\rho C_xS = \qty{2.96}{km}$, $d = \lambda\ln2 =
\qty{2.0}{km}$ (both together: $\qty{1.0}{km}$).

\textbf{6.} Weight, \omterm{def:b2:rigid-body-mechanics:contact}{normal reactions}, friction of the road on the tyres.
The forward friction on the front tyres is the only horizontal external
force, so it accelerates the car; its power is zero (contact point at
rest); the kinetic energy comes from the engine --- internal work.

\textbf{7.} $M\dot v = T_f - T_r$; front wheels $2J\dot v/R = \Gamma - T_fR$;
rear wheels $2J\dot v/R = T_rR$. Adding: $(M + 4J/R^2)\dot v = \Gamma/R$;
$M_{\text{eff}} = \qty{1242}{kg}$.

\textbf{8.} Moments about $G$ (horizontal ground forces at depth $h$,
total $M\dot v$): $aN_f - bN_r + hM\dot v = 0$ with $N_f + N_r = Mg$, so
$N_r = Mga/L + Mh\dot v/L$: the rear gains $Mh\dot v/L$, the front loses it.

\textbf{9.} Front drive: $M\dot v \le fN_f = f(Mgb/L - Mh\dot v/L)$, so
$\dot v_{\max} = fgb/(L + fh) = 11.77/3.04 = \qty{3.9}{m/s^2}$. Rear drive:
$M\dot v \le f(Mga/L + Mh\dot v/L)$, $\dot v_{\max} = fga/(L - fh) = \qty{4.0}{m/s^2}$.
Load transfer helps the rear axle: rear drive, and weight at the back,
accelerate harder.

\textbf{10.} $\Gamma = RM_{\text{eff}}\dot v_{\max} = 0.31 \times 1242 \times 3.87 =
\qty{1.5}{kN.m}$; $P = \Gamma v/R = 1490 \times 13.9/0.31 = \qty{67}{kW}$.

\textbf{11.} $\dot v_{\max} = 0.1 \times 9.81 \times 1.5/2.655 = \qty{0.55}{m/s^2}$;
\qty{25}{s} to \qty{50}{km/h}.

\textbf{12.} Gripping: \omterm{thm:b2:rigid-body-mechanics:coulomb}{static friction}, no slip, no dissipation at the
contact. Spinning: \omterm{thm:b2:rigid-body-mechanics:coulomb}{kinetic friction} (slightly smaller coefficient),
its work heats the tyre and polishes or melts the ice, and the engine's
power goes into spinning the wheel, not moving the car.

\textbf{13.} The brake torque is internal (wheel--body); the car is
slowed only by the road's backward friction on the tyre. For the wheel
(negligible $J\dot\omega$): $\Gamma_b \approx T_iR \le fN_iR$.

\textbf{14.} Same balance with $\dot v = -\gamma$: $N_f = Mgb/L + Mh\gamma/L$;
at $\gamma = fg = \qty{7.85}{m/s^2}$: $N_f = 6792 + 1992 = \qty{8.8}{kN}$, $N_r =
\qty{3.0}{kN}$.

\textbf{15.} $\gamma_{\max} = fg = \qty{7.85}{m/s^2}$; $d = v^2/2\gamma = \qty{40}{m}$;
wet $f = 0.45$: $\qty{4.4}{m/s^2}$, \qty{71}{m}.

\textbf{16.} Equal torques give equal forces, but the rear axle carries
only \qty{3.0}{kN}: it locks first. A sliding tyre's friction is opposite
to its \omterm{def:b2:rigid-body-mechanics:rolling}{slip velocity} and can no longer supply a lateral force: a locked
rear axle loses directional stability and the car spins.

\textbf{17.} Locked: $\gamma = 0.7g = \qty{6.9}{m/s^2}$, $12\%$ less, stopping
distance \qty{45}{m} instead of \qty{40}{m} --- and no steering. An
anti-lock system senses a wheel decelerating toward lock, releases and
reapplies the brake many times a second, keeping the tyre near the
static peak and rolling.

\textbf{18.} \qty{388}{kJ} into $4 \times 6 \times 470 = \qty{11.3}{kJ/K}$: $\Delta T =
\qty{34}{K}$ per stop. A \qty{40}{t} truck descending \qty{1000}{m} must dissipate
$Mgh = \qty{390}{MJ}$: brakes overheat and fade; a gravel ramp stops the
truck by the friction of the gravel.

\textbf{19.} $0.7 \times 388 = \qty{272}{kJ} = \qty{0.075}{kWh}$; a \qty{50}{kWh}
battery is about $660$ such stops.

\textbf{20.} $Mv^2/\rho \le fMg$: $v \le \sqrt{fg\rho} = \qty{19.8}{m/s} =
\qty{71}{km/h}$; wet \qty{53}{km/h}.

\textbf{21.} Lateral friction $Mv^2/\rho$ at depth $h$, normal forces at
$\pm w/2$: $(N_{\text{out}} - N_{\text{in}})\,w/2 = Mv^2h/\rho$, so $\Delta N =
Mv^2h/\rho w$. Inner wheels lift at $\Delta N = Mg/2$: $v = \sqrt{g\rho w/2h} =
\qty{25.9}{m/s} = \qty{93}{km/h} > \qty{71}{km/h}$: it skids first (since $f <
w/2h = 1.36$).

\textbf{22.} On a bank of angle $\beta$ with no friction, $N\sin\beta =
Mv^2/\rho$ and $N\cos\beta = Mg$: $v = \sqrt{g\rho\tan\beta} = \qty{9.3}{m/s} =
\qty{33}{km/h}$ for \qty{10}{\degree}.

\textbf{23.} $\omega_0 = \qty{6.28}{rad/s}$, $J\dot\omega = \qty{-0.05}{N.m}$: $t =
\qty{126}{s}$; $\tfrac12J\omega_0^2 = \qty{20}{J}$ lost.

\textbf{24.} \omterm{thm:b2:rigid-body-mechanics:coulomb}{Kinetic friction} $fN = \qty{2.4}{kN}$ pushes the car forward:
$\dot v = 2400/1200 = \qty{2.0}{m/s^2}$; on the wheel $J\dot\omega = -fNR$,
$\dot\omega = \qty{-744}{rad/s^2}$. The slip $R\omega - v$ starts at \qty{1.95}{m/s}
and falls at $\qty{233}{m/s^2}$: it vanishes in \qty{8.4}{ms}, the car having
gained \qty{1.7}{cm/s} and moved \qty{70}{\micro m} --- a jolt, not a push.

\textbf{25.} Rolling: friction $\approx 0$, rolling resistance $\mu_rN$,
$|T| \le fN$ trivially. Accelerating: forward friction on the driven
tyres, $T \le fN_{\text{driven}}$, $\dot v \le fgb/(L + fh)$. Braking:
backward friction on all tyres, $T_i \le fN_i$, $\gamma \le fg$.
Cornering: lateral friction $Mv^2/\rho \le fMg$. In every case one
force --- the road's friction, bounded by Coulomb's law --- does the
whole job.
\end{solution}
